\chapter{Conclusions and future work}
\label{cap:conclusiones}

This chapter is dedicated to gathering the conclusions drawn throughout the project and exploring the future lines of work available. 

\section{Conclusions}

Through this project, the feasibility of using CNNs to automate parameter estimation in FM synthesis has been demonstrated. The developed system succeeds in reducing the technical barrier of \textit{sound matching}, confirming that deep learning is an effective tool to streamline the workflow of producers and sound designers, allowing them to focus on the creative process. 

This has been corroborated by comparing it with traditional models such as \textit{grid search}, thus demonstrating its superiority in efficiency and computational viability (see Section \ref{sec:gridSearch}).

Additionally, with the experimentation carried out for Chapter \ref{cap:ResultadosYEvaluacion}, different architecture and configuration options have been contrasted, obtaining unexpected results from which it could be concluded that complex architectures do not guarantee an increase in performance, as they require extensive training and a large \textit{dataset} size.

Furthermore, it has been shown that mel spectrograms perform better than linear ones both in perceptual terms, by yielding better results, and in computational terms, since a smaller number of \textit{bins} is required to generate the spectrogram.

Lastly, it is worth highlighting the importance that the application of graphical interfaces and \textit{pipeline} automation has had in the project, through which it has been possible to easily manage different architectures, training sessions, and tests. 

\section{Future work}

A work of this nature is easily continuable from multiple perspectives; however, we consider it relevant to highlight the four most promising branches of research and development:

\begin{enumerate}
\item \textbf{Implementation of a differentiable synthesizer:} Shifting the project's focus towards the use and adaptation of DDSP (\textit{Differentiable Digital Signal Processing}), which would allow the synthesis process to be integrated directly into the network's training. This would enable exact calculation in gradient descent (by not treating the synthesizer as a ``black box''), optimizing the model's convergence and significantly improving the fidelity of the results.

\item \textbf{Scalability to other synthesis methods:} The established workflow lays the foundation for expanding the system beyond frequency modulation synthesis. It would be of great interest to adapt the model to estimate parameters in additive, subtractive, physical modeling, or granular synthesis.

\item \textbf{Architectural exploration and optimization:} To overcome current performance limits, a more exhaustive study of hyperparameters, the evaluation of alternative network architectures, and the design of custom loss functions are proposed. This would allow for the contribution of novel and proprietary approaches, reducing the reliance on standardized models in current literature.

\item \textbf{Development of a VST plugin:} Looking ahead to its commercial and practical viability, the logical step is to encapsulate the model in a VST format for direct integration into any DAW (\textit{Digital Audio Workstation}). This will require the design of an intuitive graphical user interface (GUI) and the programming of a more robust system for exception handling, preventing the application from crashing due to unforeseen inputs or \textit{edge cases} from the end user.

\item \textbf{Perceptual evaluation through listening tests:} Although the model's performance has been measured using quantitative metrics, incorporating qualitative validation through real users would represent a significant advancement. Conducting subjective perception tests with producers and sound designers would allow for the evaluation of timbral similarity from a real psychoacoustic perspective, providing an empirical success metric directly aligned with the human ear.
\end{enumerate}